\chapter{Methodisches Vorgehen und Anforderungen}
\label{cha:methodik}

% Wie wurde gearbeitet (Vorgehensmodell, Erhebung, Werkzeuge) und was muss die
% Lösung leisten? Anforderungen nummerieren, damit sie in der Evaluation
% wieder aufgegriffen werden können.

\section{Vorgehensmodell}
\label{sec:vorgehen}

Die Arbeit folgt einem iterativen Vorgehen aus Anforderungsanalyse, Entwurf,
Implementierung und Evaluation (\vgl Abbildung~\ref{fig:vorgehen}).

\begin{figure}[htbp]
  \centering
  % \includegraphics[width=0.8\textwidth]{vorgehen}
  \fbox{\parbox{0.78\textwidth}{\centering\vspace{1.4cm}Platzhalter: Vorgehensmodell als Grafik (abbildungen/vorgehen.png)\vspace{1.4cm}}}
  \caption{Iteratives Vorgehen der Arbeit}
  \label{fig:vorgehen}
\end{figure}

\section{Anforderungen}
\label{sec:anforderungen}

% Funktionale Anforderungen (was das System tun soll) und nicht-funktionale
% Anforderungen (wie gut: Leistung, Usability, Sicherheit …).
% Priorität: M = Muss, S = Soll, K = Kann.

\begin{table}[htbp]
  \centering
  \caption{Funktionale Anforderungen}
  \label{tab:anforderungen-funktional}
  \begin{tabularx}{\textwidth}{lXc}
    \toprule
    Nr. & Anforderung & Prio \\
    \midrule
    FA1 & Ein bestehender Projektordner lässt sich öffnen und durchsuchen. & M \\
    FA2 & Änderungen werden automatisch kompiliert und als PDF angezeigt. & M \\
    FA3 & Jede Version bleibt über eine Historie wiederherstellbar. & S \\
    FA4 & Literaturangaben lassen sich per DOI importieren. & K \\
    \bottomrule
  \end{tabularx}
\end{table}

\begin{table}[htbp]
  \centering
  \caption{Nicht-funktionale Anforderungen}
  \label{tab:anforderungen-nichtfunktional}
  \begin{tabularx}{\textwidth}{lXc}
    \toprule
    Nr. & Anforderung & Prio \\
    \midrule
    NFA1 & Die Kompilierung unterbricht den Schreibfluss nicht (asynchron). & M \\
    NFA2 & Alle Daten bleiben lokal auf dem Rechner. & M \\
    NFA3 & Die Anwendung startet in unter fünf Sekunden. & S \\
    \bottomrule
  \end{tabularx}
\end{table}

\section{Bewertungskriterien}
\label{sec:kriterien}

% Woran wird der Erfolg gemessen? Diese Kriterien strukturieren Kapitel 5.
Die Lösung wird anhand folgender Kriterien bewertet:
\begin{description}
  \item[Funktionalität] Erfüllungsgrad der Anforderungen aus Tabelle~\ref{tab:anforderungen-funktional}.
  \item[Leistung] Dauer einer Kompilierung und Reaktionszeit der Oberfläche.
  \item[Nutzbarkeit] Ergebnisse einer Nutzerbefragung (\zB System Usability Scale).
\end{description}
